\begin{abstract}
Parameter-efficient adaptation of MedSAM to one imaging domain can leave it worse than zero-shot on distant domains, a collapse that earlier work linked to drift of the mask decoder's representations from the base model. We intervene on that mechanism: a centred kernel alignment (CKA) loss pulls the adapted decoder's activations on a 32-image out-of-distribution probe back toward the base decoder during LoRA training on dermoscopy. Over three seeds the regulariser raises tight-box Dice by 0.05 on breast ultrasound and 0.10 on mammography with no in-domain cost; the same anchor on the image encoder alone lowers far-OOD Dice, and anchoring encoder and decoder together gives the largest gain. Varying the adaptation budget from 50 to 2595 images shows that LoRA falls below zero-shot far-OOD already at 250 images, whereas decoder-only and encoder-only adaptation never do. Finally, the per-image drift of the decoder's upscaled mask embedding detects far-OOD failures of parameter-efficient adapters with AUROC 0.95 on ultrasound and 0.82 on mammography, where the decoder's own IoU estimate is at chance, while the IoU estimate is informative for full fine-tuning. Far-OOD robustness of adapted MedSAM is decided at the decoder, and it can be both protected and monitored there.
\keywords{Medical image segmentation \and Foundation models \and Parameter-efficient fine-tuning \and Out-of-distribution robustness \and Failure detection}
\end{abstract}
